% ============================================================
%  Marcadores y inclusiones "seguras"
%  Objetivo: que Overleaf compile aunque falte una figura o tabla,
%  y que lo que falta quede MUY visible en el PDF.
% ============================================================

% --- Hueco reservado para una imagen (de la plantilla original) ---
% Uso:   \imgplaceholder{texto descriptivo}
%        \imgplaceholder[alto]{texto descriptivo}   (alto por defecto: 5cm)
\newcommand{\imgplaceholder}[2][5cm]{%
  \begin{center}
    \fboxsep=8pt
    \textcolor{gray}{\fbox{\begin{minipage}[c][#1][c]{0.82\linewidth}
      \centering\itshape\small #2
    \end{minipage}}}
  \end{center}
}

% --- Incluir figura; si el archivo no existe, muestra un hueco ---
% Uso: \safegraphics[opciones]{figures/nombre.pdf}{descripción si falta}
% (la ruta debe llevar extensión, porque \IfFileExists la necesita)
\newcommand{\safegraphics}[3][width=\linewidth]{%
  \IfFileExists{#2}%
    {\includegraphics[#1]{#2}}%
    {\imgplaceholder[3cm]{Falta la figura: \texttt{\detokenize{#2}}\\ #3}}%
}

% --- Incluir tabla generada (solo el entorno tabular, sin flotante) ---
% Uso: \safetable{generated/tables/phase_summary}   (sin .tex)
\newcommand{\safetable}[1]{%
  \IfFileExists{#1.tex}%
    {\input{#1}}%
    {\textcolor{orange}{\textbf{[Falta la tabla: \texttt{\detokenize{#1}}]}}}%
}

% --- Contenido pendiente de redactar (visible en naranja) ---
\newcommand{\todoslot}[1]{\textcolor{orange}{\textbf{[PENDIENTE: #1]}}}
